% commentaire.tex — f170v, vv. 679–720 : commentaire de la pièce (conversion des travaux antérieurs ;
% discussions textuelles passées dans critique.tex, voir repartition.md)
\comm{680}{μανίαις φλέγων λύσσῃ παράκοπος.}{Παράκοπος, \italique{délirant}, est une métaphore usée\nc{Eschyle \italique{Prom.} 581 ; Eur. \italique{Bacch.} 33, 1000, etc.} et d'origine incertaine. Elle paraît sans rapport avec la frappe de la monnaie\nc{Malgré παρακόπτειν (\italique{Ach.} 517), « battre de la fausse monnaie ».} ; mais elle en a peut-être un avec la notion d'« ébranler le cerveau »\nc{Cf. §474 du même livre : \italique{Les images d'Aristophane} de Jean Taillardat.}, cf. Euripide \italique{Hippolyte}, v. 238, etc... παρακόπτειν φρένας. (Note de Taillardat §476.)}
\comm{683}{γυναιξὶ καὶ βροτοῖς.}{L'on s'attendait à la formule consacrée : θεοῖσι καὶ βροτοῖς\nc{Homère, \italique{Iliade} XXIV, 553 ; Pindare, \italique{Pythiques} 3.30 ; Sophocle, \italique{Électre} v. 199.}.}
\comm{689–761}{Parodie du \italique{Télèphe} d'Euripide.}{Comme dans ses \italique{Acharniens} représentés une quinzaine d'années plus tôt -- il y a de nombreuses analogies entre ces deux pièces --, Aristophane détourne le \italique{Télèphe} d'Euripide. On ne sait pas grand'chose de cette dernière pièce, la majorité des fragments conservés étant eux-mêmes tirés d'Aristophane et de ses anciens commentateurs\nc{Pour une présentation détaillée du \italique{Télèphe}, le lecteur pourra consulter l'édition de François Jouan \& Herman van Looy aux Belles Lettres, tome VIII, \nieme{3} partie.}.
Télèphe était un héros populaire puisque les trois grands tragiques lui ont consacré pas moins de sept pièces ; il fut blessé par un coup de lance d'Achille.
Comme sa blessure ne guérissait pas, il finit par consulter l'oracle d'Apollon, réponse de ce dernier : ὁ τρώσας ἰάσεται, « Qui t'a blessé te guérira ». Il se rend donc, déguisé en mendiant, dans le camp de l'armée achéenne préparant son expédition sur Troie. Allait-il d'abord trouver Clytemnestre qui lui confiait son petit Oreste pour appuyer son plaidoyer ? toujours est-il que Télèphe se retrouva la proie exposée à la colère des chefs grecs, il se réfugiait alors en suppliant sur l'autel en prenant en otage Oreste. Le sort de l'enfant était plus heureux que dans sa thesmophorienne parodie : grâce à l'entremise d'Agamemnon, Achille se laissait convaincre de soigner Télèphe et tout était bien qui finissait bien. Outre cette scène d'enlèvement du petit Oreste, de nombreux autres passages des \italique{Thesmophories} sont sans doute aussi inspirés du \italique{Télèphe} (v. 39-70 : où l'altercation entre le héros et un gardien attirait Clytemnestre à l'extérieur et permettait sa rencontre avec Télèphe ; v. 466-470 : lorsqu'au début de son plaidoyer, le héros -- sous une fausse identité -- usait de précautions oratoires ; v. 598-600, 659-667 avec l'annonce de la présence d'un espion dans le camp grec.)}
\comm{693–5}{ἐπὶ τῶν μηρίων ... βωμόν.}{Vers tragiques, très sûrement empruntés à Euripide. (Note de Willems.)}
\comm{696–7}{οὐ πολλὴν βοὴν / στήσεσθε καὶ τροπαῖον.}{Dans le sens de μήστωρ φόβοιο d'Homère. « Un vacarme à faire sauver ». (Musset, \italique{Confess}., p. 111 ; note de Willems.)}
\comm{702}{ὡς ἅπαν γὰρ ἐστὶν τόλμης ἔργα κἀναισχυντίας.}{Le scholiaste met en avant le tour proverbial de la formule en commentant d'un sobre «~παροιμία~». Heureusement, d'autres glosateurs furent davantage prolixes. On retrouve la formule chez Suidas (α 2902) : 
\citComm{ἅπαντα τόλμης πλέα κἀναισχυντίας: ἐπὶ τῶν ἰταμῶν καὶ ἀναιδῶν.}
Et dans le \italique{Corpus paroemiographorum Graecorum}\nc{\italique{Corpus paroemiographorum Graecorum}, vol. I, Appendices, I, 35, p. 384 ($=$ Vat. K II 36).} :
\citComm{ἅπαντα τόλμης καὶ ἀναισχυντίας: ὅτι ἡ ἀναισχυντία νικᾷ τὰ πολλά.}
Enfin dans le codex \italique{Laurentianus} LVIII 24\nc{Il s'agit d'un manuscrit conservé à la Biblioteca Medicea Laurenziana (BML) de Florence, dépôt Plutarque, sous la cote 58. 24 et édité par Leopold Cohn en 1887 : \italique{Zu den Paroemiographen} (notre proverbe y figure à la page 40).}, l'on tombe sur la formulation suivante :\citComm{ἅπαντα τόλμης δοῦλα καὶ ἀναισχυντίας.}Ces trois leçons sont évidemment tirées du vers 702 des \italique{Thesmophories} et des anciens commentaires afférents. Pour le texte, voir le commentaire critique.

Finissons cette note aporétique en remarquant que les termes τόλμη et ἀναισχυντία sont souvent mis en collocation : Platon\nc{Platon, \italique{Apologie de Socrate}, 38d. Voir aussi son \italique{Théétète} (196d) où le tour proverbial est encore plus sensible : ὅμως δέ — πάντα γὰρ τολμητέον — τί εἰ ἐπιχειρήσαιμεν ἀναισχυντεῖν; «~De l'audace ! encore de l'audace ! toujours de l'audace !~»} et Démosthène\nc{Παραγραφὴ πρὸς Ναυσίμακον καὶ Ξενοπείθην, 5 (Νόμος).} associent les deux mots avec un sens analogue. Isocrate nous apporte quelque précision, l'on comprend en le lisant\nc{Περὶ ἀντιδόσεως, 190 : ἔτι δὲ τὴν τόλμαν μὴ τὴν ἀναισχυντίας σημεῖον γιγνομένην, ἀλλὰ τὴν μετὰ σωφροσύνης. Voir aussi Démosthène, \italique{In Aristogitonem 1}, 68.8 : ταῦτα τολμᾷ λέγειν δι' ἀναισχυντίαν.} que l'audace (τόλμα) pouvait être interprétée comme un signe (σημεῖον) d'impudence (ἀναισχυντία). Les Grecs avaient fait d'\italique{Implacabilité} un dieu : θεὸς ἡ Ἀναίδεια\nc{\italique{Implacabilité} et non \italique{Impudence}. Cf. le \italique{Dictionnaire des antiquités grecques et romaines} (Daremberg \& Saglio), t. I, p. 398b (\latin{Areopagus}). Les Athéniens avaient dressé des autels à Ὕβρις et Ἀναίδεια dans l'Aréopage, contre lesquels s'était insurgé ce rabat-joie de Cicéron -- non sans faire un contre-sens comme le souligne l'auteur de l'article. Dans sa querelle avec Agamemnon, l'attitude d'Achille -- qui par bien des aspects est celle d'un adolescent attardé -- nous révèle en fait son côté divin : il est \italique{implacable}, tout autant que le destin qu'il connaît et a accepté. Pour revenir au proverbe `θεὸς ἡ Ἀναίδεια', certainement inspiré par ces autels de l'Aréopage, voir Zénobius (IV, 36). Érasme l'a aussi recensé dans ses \italique{Adages} (n° 2063).}.}
\comm{704}{οἷον ὑμῶν ἐξάρξω τὴν ἄγαν αὐθαδίαν.}{Pour le texte, voir le commentaire critique. L'occasion fait le larron : profitons-en pour placer cet aphorisme rapporté par Athénée\nc{Athénée, II, 37b.} illustrant le sens de la mesure grec, même quand il s'agit de boire...\citComm{Πάσας δ᾽ ἐκ κραδίας ἀνίας ἀνδρῶν ἀλαπάζει\\
πινόμενος κατὰ μέτρον· ὑπὲρ μέτρον δὲ χερείων.\\
\francais{Le vin est souverain pour chasser les soucis du cœur de qui n'en boit pas trop, mais l'excès nuit.}}(Un franc-buveur espiègle répondrait que la mesure du boire est de boire sans mesure\nc{Aristophane, \italique{Guêpes}, vers 79-80 :\citComm{εἶναι φιλοπότην αὐτόν. [ ... ]\\
αὕτη γε χρηστῶν ἐστιν ἀνδρῶν ἡ νόσος.}.} !..)}
\comm{710–2}{ἀλλ᾽ οὖν ἥκεις τ᾿ ὅθεν ἥκεις.}{Le chœur enjoint Mnésiloque de cesser cette escalade de l'impudence : qu'il quitte au moins l'autel et les rejoigne\nc{C'est l'explication de la scholie aussi conservée par Suidas (η 170) : ἥκεις φησίν εἰς τὰς χεῖρας ἡμῶν, ὅθεν οὐ φεύξῃ.}. Pour le texte, voir le commentaire critique.}
